\section{Actual whole original pulled projective lines and their complete chart coefficients}
\label{sec:native-pulled-projective-line-frames}

\begin{theorem}[Actual whole pulled lines retain every original chart frame {\cite[\S1.2]{BJ22}}]
\label{thm:native-pulled-projective-line-frames}
\lean{EnvelopeProperty.NativePulledProjectiveLineFrames.line, EnvelopeProperty.NativePulledProjectiveLineFrames.line_isInvertible, EnvelopeProperty.NativePulledProjectiveLineFrames.chartFrame, EnvelopeProperty.NativePulledProjectiveLineFrames.chartCoefficient, EnvelopeProperty.NativePulledProjectiveLineFrames.formSection, EnvelopeProperty.NativePulledProjectiveLineFrames.chartCoefficient_formSection}
\leanok
\uses{lem:fc-original-homogeneous-form-frame-coefficient,thm:fc-original-section-pullback-model-norm}
An actual original closed projective embedding gives its same whole
pulled line $f^*O(d)$, its genuine restriction frame on every original
coordinate inverse-image chart, and the actual original regular coefficient
of every whole original global section. For every original homogeneous form
$F$ of degree $d$, its genuine pulled global section has exactly the original
pulled coefficient $F/X_j^d$ on every entire original chart. No frame or
section coefficient comparison is an input.
\end{theorem}
\begin{proof}
Pull the actual original projective frame back to the actual original
inverse-image chart. Native pullback composition and the original structure
unit comparison give a whole line trivialization; the genuine restriction
versus open-pullback equivalence gives its genuine restriction frame. Its
actual original section-module equivalence computes the coefficient of
every original global section. The original global homogeneous section is
pulled as its whole structure-line morphism, evaluated at the original unit.
The genuine whole local scalar equation for the original homogeneous form
pulls back to the actual inverse-image chart. The genuine original local
coefficient equation for the resulting global section has the same whole
left-hand morphism. Evaluation at the actual unit and injectivity of the
original top-section open isomorphism identify the two whole original
regular coefficients.
\end{proof}

\begin{theorem}[Actual generic homogeneous coefficients and nonzero pivot sections {\cite[\S1.2]{BJ22}}]
\label{thm:native-pulled-projective-line-pivot}
\lean{EnvelopeProperty.NativePulledProjectiveLineFrames.germ_chartCoefficient_formSection, EnvelopeProperty.NativePulledProjectiveLineFrames.coordinatePowerSection, EnvelopeProperty.NativePulledProjectiveLineFrames.germ_chartCoefficient_coordinatePowerSection, EnvelopeProperty.NativePulledProjectiveLineFrames.coordinatePowerSection_ne_zero}
\leanok
\uses{thm:native-pulled-projective-line-frames,thm:projective-native-homogeneous-coefficient}
Assume throughout this statement that the original scheme $X$ is integral.
On every original nonempty coordinate chart, the actual generic germ of the
coefficient of an original homogeneous section is exactly its original
polynomial evaluated at the actual native coordinate ratios. In particular
the original section given by the power of that pivot coordinate is a
nonzero whole original global section of the pulled line, internally
normalized by coefficient one in the actual original function field.
\end{theorem}
\begin{proof}
Apply the preceding entire coefficient identity and the genuine native
homogeneous generic-coefficient theorem. For the pivot power, polynomial
evaluation is the power of the actual pivot ratio, which is one. A zero
whole global section would have zero actual original chart coefficient and
zero native generic germ, contradicting this literal field value one. This
works also in degree zero and uses no supplied nonzero-section certificate.
\end{proof}
